\begin{thebibliography}{99}
\bibitem{AlekseevBradford2026} V.~Alekseev and H.~Bradford, Sofic actions, halo products, and metric approximations of groups, \href{https://arxiv.org/abs/2601.18742}{arXiv:2601.18742} (2026).
\bibitem{AD2006} C.~Anantharaman-Delaroche, On ergodic theorems for free group actions on noncommutative spaces, Probab. Theory Related Fields \textbf{135} (2006), 520--546.
\bibitem{ArakiYamagami1981} H.~Araki and S.~Yamagami, An inequality for Hilbert--Schmidt norm, Comm. Math. Phys. \textbf{81} (1981), 89--96.
\bibitem{BhatDevendra2025} B.~V.~R.~Bhat and R.~Devendra, On regions of mixed unitarity for semigroups of unital quantum channels, \href{https://arxiv.org/abs/2512.23598}{arXiv:2512.23598} (2025).
\bibitem{Bradford2022} H.~Bradford, Controlling LEF growth in some group extensions, \href{https://arxiv.org/abs/2201.05052}{arXiv:2201.05052} (2022).
\bibitem{BHV2006} S.~Bravyi, M.~B.~Hastings and F.~Verstraete, Lieb--Robinson bounds and the generation of correlations and topological quantum order, Phys. Rev. Lett. \textbf{97} (2006), 050401.
\bibitem{Bridson2002} M.~R.~Bridson, The geometry of the word problem, in: M.~R.~Bridson and S.~M.~Salamon (eds.), Invitations to Geometry and Topology, Oxf. Grad. Texts Math. \textbf{7}, Oxford University Press, Oxford, 2002, 29--91.
\bibitem{Brin2004} M.~G.~Brin, Higher dimensional Thompson groups, Geom. Dedicata \textbf{108} (2004), 163--192.
\bibitem{Brin2005} M.~G.~Brin, Presentations of higher dimensional Thompson groups, J. Algebra \textbf{284} (2005), 520--558; arXiv:math/0501082.
\bibitem{Brin2010} M.~G.~Brin, On the baker's map and the simplicity of the higher dimensional Thompson groups $nV$, Publ. Mat. \textbf{54} (2010), 433--439; arXiv:0904.2624.
\bibitem{Brown1987} K.~S.~Brown, Finiteness properties of groups, J. Pure Appl. Algebra \textbf{44} (1987), 45--75.
\bibitem{BCMR2016} J.~Burillo, S.~Cleary, A.~Martino and C.~E.~R\"over, Commensurations and metric properties of Houghton's groups, Pacific J. Math. \textbf{285} (2016), 289--301; arXiv:1403.0026.
\bibitem{BCS2001} J.~Burillo, S.~Cleary and M.~I.~Stein, Metrics and embeddings of generalizations of Thompson's group $F$, Trans. Amer. Math. Soc. \textbf{353} (2001), 1677--1689; arXiv:math/9809185.
\bibitem{CFP1996} J.~W.~Cannon, W.~J.~Floyd and W.~R.~Parry, Introductory notes on Richard Thompson's groups, Enseign. Math. (2) \textbf{42} (1996), 215--256.
\bibitem{CapraroLupini2015} V.~Capraro and M.~Lupini, Introduction to Sofic and Hyperlinear Groups and Connes' Embedding Conjecture, Lecture Notes in Math. \textbf{2136}, Springer, Cham, 2015.
\bibitem{CarterWegman1979} J.~L.~Carter and M.~N.~Wegman, Universal classes of hash functions, J. Comput. System Sci. \textbf{18} (1979), 143--154.
\bibitem{Chiribella2009} G.~Chiribella, G.~M.~D'Ariano and P.~Perinotti, Theoretical framework for quantum networks, Phys. Rev. A \textbf{80} (2009), 022339; arXiv:0904.4483.
\bibitem{Chou1980} C.~Chou, Elementary amenable groups, Illinois J. Math. \textbf{24} (1980), 396--407.
\bibitem{Cornulier2006} Y.~de Cornulier, Finitely presented wreath products and double coset decompositions, Geom. Dedicata \textbf{122} (2006), 89--108; arXiv:math/0509090.
\bibitem{Cornulier2013} Y.~Cornulier, Sofic profile and computability of Cremona groups, Michigan Math. J. \textbf{62}(4) (2013), 823--841; \href{https://doi.org/10.1307/mmj/1387226167}{doi:10.1307/mmj/1387226167}. Numbered references use \href{https://arxiv.org/abs/1305.0993v1}{arXiv:1305.0993v1}.
\bibitem{DouglasMghirbiA} S.~Douglas and N.~Mghirbi, Houghton's group $H_3$ has superpolynomial sofic profile, version 1.1.2, Zenodo (2026); \href{https://doi.org/10.5281/zenodo.23070869}{doi:10.5281/zenodo.23070869}. Section numbers refer to version 1.1.2.
\bibitem{DouglasMghirbiB} S.~Douglas and N.~Mghirbi, Causal quantum-channel simulation: memory beyond entropy, version 1.1.2, Zenodo (2026); \href{https://doi.org/10.5281/zenodo.23070870}{doi:10.5281/zenodo.23070870}. Section numbers refer to version 1.1.2.
\bibitem{GKEP2025} D.~Gao, S.~Kunnawalkam Elayavalli and G.~Patchell, Soficity for group actions on sets and applications, Res. Math. Sci. \textbf{12} (2025), 48; arXiv:2401.04945.
\bibitem{Genevois2019} A.~Genevois, Embeddings into Thompson's groups from quasi-median geometry, Groups Geom. Dyn. \textbf{13} (2019), 1457--1510; \href{https://arxiv.org/abs/1709.03888v2}{arXiv:1709.03888v2}. Numbered references use this arXiv version.
\bibitem{GenevoisTessera2024} A.~Genevois and R.~Tessera, Lamplighter-like geometry of groups, \href{https://arxiv.org/abs/2401.13520}{arXiv:2401.13520} (2024).
\bibitem{Guba2006} V.~S.~Guba, The Dehn function of Richard Thompson's group $F$ is quadratic, Invent. Math. \textbf{163} (2006), 313--342.
\bibitem{GutoskiWatrous2007} G.~Gutoski and J.~Watrous, Toward a general theory of quantum games, Proceedings of the 39th Annual ACM Symposium on Theory of Computing (STOC 2007), 565--574; arXiv:quant-ph/0611234.
\bibitem{HaagerupMusat2011} U.~Haagerup and M.~Musat, Factorization and dilation problems for completely positive maps on von Neumann algebras, Comm. Math. Phys. \textbf{303} (2011), 555--594; arXiv:1009.0778.
\bibitem{HaagerupMusat2015} U.~Haagerup and M.~Musat, An asymptotic property of factorizable completely positive maps and the Connes embedding problem, Comm. Math. Phys. \textbf{338} (2015), 721--752; arXiv:1408.6476.
\bibitem{Harper1964} L.~H.~Harper, Optimal assignments of numbers to vertices, J. Soc. Indust. Appl. Math. \textbf{12} (1964), 131--135.
\bibitem{HayesSale2018} B.~Hayes and A.~W.~Sale, Metric approximations of wreath products, Ann. Inst. Fourier \textbf{68} (2018), 423--455; arXiv:1608.02610.
\bibitem{HennigMatucci2012} J.~Hennig and F.~Matucci, Presentations for the higher-dimensional Thompson groups $nV$, Pacific J. Math. \textbf{257} (2012), 53--74; arXiv:1105.3714.
\bibitem{HoffmanWielandt1953} A.~J.~Hoffman and H.~W.~Wielandt, The variation of the spectrum of a normal matrix, Duke Math. J. \textbf{20} (1953), 37--39.
\bibitem{HHO2003} M.~Horodecki, P.~Horodecki and J.~Oppenheim, Reversible transformations from pure to mixed states and the unique measure of information, Phys. Rev. A \textbf{67} (2003), 062104.
\bibitem{Houghton1978} C.~H.~Houghton, The first cohomology of a group with permutation module coefficients, Arch. Math. \textbf{31} (1978), 254--258.
\bibitem{Ivanov2021} A.~Ivanov, Sofic profiles of $S(\omega)$ and computability, Arch. Math. Logic \textbf{60} (2021), 477--494; arXiv:2001.02648.
\bibitem{JNVWY2020} Z.~Ji, A.~Natarajan, T.~Vidick, J.~Wright and H.~Yuen, MIP$^*$ = RE, \href{https://arxiv.org/abs/2001.04383}{arXiv:2001.04383} (2020).
\bibitem{Johnson1999} D.~L.~Johnson, Embedding some recursively presented groups, in: Groups St Andrews 1997 in Bath, II, London Math. Soc. Lecture Note Ser. \textbf{261}, Cambridge Univ. Press, Cambridge, 1999, 410--416.
\bibitem{JungeLeMerdyXu2006} M.~Junge, C.~Le Merdy and Q.~Xu, $H^\infty$ functional calculus and square functions on noncommutative $L^p$-spaces, Ast\'erisque \textbf{305} (2006), vi+138 pp.
\bibitem{Krengel1985} U.~Krengel, Ergodic Theorems, de Gruyter Stud. Math. \textbf{6}, Walter de Gruyter, Berlin, 1985.
\bibitem{KretschmannWerner2005} D.~Kretschmann and R.~F.~Werner, Quantum channels with memory, Phys. Rev. A \textbf{72} (2005), 062323; arXiv:quant-ph/0502106.
\bibitem{Kummerer1985} B.~K\"ummerer, Markov dilations on W$^*$-algebras, J. Funct. Anal. \textbf{63} (1985), 139--177.
\bibitem{KunThom2026} G.~Kun and A.~Thom, Nonsofic wreath products of residually finite groups, \href{https://arxiv.org/abs/2608.06222}{arXiv:2608.06222} (2026).
\bibitem{Lee2012} S.~R.~Lee, Geometry of Houghton's groups, PhD thesis, University of Oklahoma (2012); \href{https://arxiv.org/abs/1212.0257v1}{arXiv:1212.0257v1}. Numbered references use this arXiv version.
\bibitem{LiebRobinson1972} E.~H.~Lieb and D.~W.~Robinson, The finite group velocity of quantum spin systems, Comm. Math. Phys. \textbf{28} (1972), 251--257.
\bibitem{LyndonSchupp1977} R.~C.~Lyndon and P.~E.~Schupp, Combinatorial Group Theory, Ergeb. Math. Grenzgeb. \textbf{89}, Springer, Berlin, 1977.
\bibitem{Milnor1971} J.~Milnor, Introduction to Algebraic $K$-Theory, Ann. of Math. Stud. \textbf{72}, Princeton Univ. Press, Princeton, 1971.
\bibitem{Musat2025} M.~Musat, The Connes--Kirchberg problem and infinite-dimensional phenomena in quantum information theory, \href{https://arxiv.org/abs/2512.00432}{arXiv:2512.00432} (2025).
\bibitem{MusatRordam2020} M.~Musat and M.~R{\o}rdam, Non-closure of quantum correlation matrices and factorizable channels that require infinite dimensional ancilla (with an appendix by N.~Ozawa), Comm. Math. Phys. \textbf{375} (2020), 1761--1776; arXiv:1806.10242.
\bibitem{NSY2019} B.~Nachtergaele, R.~Sims and A.~Young, Quasi-locality bounds for quantum lattice systems. Part I. Lieb--Robinson bounds, quasi-local maps, and spectral flow automorphisms, J. Math. Phys. \textbf{60} (2019), 061101; arXiv:1810.02428.
\bibitem{Rover1999} C.~E.~R\"over, Subgroups of finitely presented simple groups, PhD thesis, University of Oxford (1999).
\bibitem{RybarZiman2008} T.~Ryb\'ar and M.~Ziman, Repeatable quantum memory channels, Phys. Rev. A \textbf{78} (2008), 052114; arXiv:0808.3851.
\bibitem{Slofstra2018} W.~Slofstra, A group with at least subexponential hyperlinear profile, \href{https://arxiv.org/abs/1806.05267v1}{arXiv:1806.05267v1} (2018).
\bibitem{VershikGordon1997} A.~M.~Vershik and E.~I.~Gordon, Groups that are locally embeddable in the class of finite groups, Algebra i Analiz \textbf{9}(1) (1997), 71--97; English transl. St. Petersburg Math. J. \textbf{9}(1) (1998), 49--67.
\end{thebibliography}
